\chapter{系统设计文档}

\providecommand{\diagramimage}[3][0.95\textwidth]{
\begin{figure}[H]
  \centering
  \includegraphics[width=#1,height=0.72\textheight,keepaspectratio]{figures/#2}
  \caption{#3}
\end{figure}
}
\providecommand{\flowdiagramimage}[3][0.62\textwidth]{
\begin{figure}[H]
  \centering
  \includegraphics[width=#1,height=0.72\textheight,keepaspectratio]{figures/#2}
  \caption{#3}
\end{figure}
}
\providecommand{\systemscreenimage}[3][0.88\textwidth]{
\begin{figure}[H]
  \centering
  \includegraphics[width=#1,height=0.72\textheight,keepaspectratio]{figures/screenshots/#2}
  \caption{#3}
\end{figure}
}

\section{引言}

\subsection{设计目标}

本文档在需求分析基础上，对语界 LinguaSpace 多语文旅智能导览与人才培养系统进行系统设计说明。文档分为概要设计和详细设计两个部分：概要设计说明系统总体架构、模块划分、部署结构和关键设计原则；详细设计进一步说明各子系统类结构、接口设计、数据设计、状态流转、异常处理和运维设计。

系统设计的核心目标是将游客多语导览、拍照识别、AI 导游实训、真人导游协同、知识库管理和内容审核统一组织为可维护的软件系统。系统不允许大语言模型直接自由生成文化事实，而是通过 RAG 文旅知识库、文化知识图谱、术语表、审核机制和日志链路保证回答可追踪、可审核、可迭代。

\subsection{设计依据}

设计依据来自系统已确定的功能边界、非功能约束、用例模型、数据模型、类模型和业务流程。技术选型围绕多语文旅导览、可信知识检索、AI 实训评分、真人导游协同和管理端审核展开，后端采用 FastAPI，前端采用 React 18、TypeScript 与 Vite，数据层以 MySQL 为运行时主库，并接入 Redis、MinIO、Neo4j；PostgreSQL/pgvector 作为后续向量检索增强预留。AI 能力层接入 Ollama 文本模型、Ollama 视觉模型、faster-whisper ASR、术语表翻译适配器和 Windows SAPI TTS。当前通过 Docker Compose 与启动脚本完成本地一键部署。

\subsection{设计原则}

\begin{enumerate}
  \item 高内聚低耦合：游客导览、学生实训、导游协同、知识管理、审核和日志模块职责明确，模块之间通过接口协作。
  \item 知识优先：历史文化、民族习俗、宗教礼仪、景区政策、开放时间和票价等内容必须优先检索审核通过的知识库。
  \item 模型可替换：语音识别、翻译、视觉识别、大语言模型、向量化和语音合成均通过适配器封装，避免业务代码绑定具体模型。
  \item 全链路可观测：每次 AI 调用记录请求编号、输入类型、检索来源、模型版本、生成结果、耗时和异常信息。
  \item 可部署可运维：系统采用容器化部署、服务分层、统一配置和监控告警机制，保证上线运行、扩容迁移和故障排查具备明确支撑。
\end{enumerate}

\subsection{系统运行能力}

系统运行能力覆盖前端交互、后端接口、模型编排、知识检索、数据持久化和运维监控。后端使用 FastAPI 暴露 REST API，前端通过环境变量配置 API 地址，模型通过统一 Provider 接入，数据库侧使用 MySQL 保存业务数据和知识切片，并通过应用层检索、Redis 缓存、MinIO 对象存储与 Neo4j 图谱镜像协同运行，RAG 链路返回可追踪来源列表。系统以完整业务闭环为目标，支持景区导览、学生实训、导游协同、知识审核和日志追踪的连续运行。

\begin{longtable}{p{3.5cm}p{4.5cm}p{5.1cm}}
\toprule
运行能力 & 接口/模块 & 设计约束 \\
\midrule
健康检查 & \code{GET /api/health} & 检查 API、PostgreSQL/pgvector、Redis、MinIO、Neo4j、LLM、Vision、TTS 和 ASR 状态。 \\
架构审计 & \code{GET /api/architecture/audit} & 对照五层架构返回达成状态，可作为管理端系统状态页面的数据来源。 \\
RAG 问答 & \code{POST /api/chat} & 先检索知识库，命中为空时返回“暂无可靠资料”，响应包含 sources、provider、model、reliable。 \\
知识库管理 & \code{GET/POST /api/knowledge}、\code{POST /api/knowledge/search} & 支持管理端查看、添加和检索已审核知识。 \\
图片问答 & \code{POST /api/image/ask} & 使用视觉模型提取图片摘要，再进入 RAG 问答链路。 \\
语音与 TTS & \code{POST /api/audio/*}、\code{POST /api/tts/synthesize} & ASR 由 faster-whisper 完成，TTS 使用 Windows SAPI 生成可播放语音文件。 \\
图谱与路线 & \code{GET /api/graph}、\code{POST /api/graph/query}、\code{POST /api/route/recommend} & 支持文化关系展示、实体关系查询和路线推荐。 \\
实训与协同 & \code{POST /api/training/score}、\code{POST /api/collaboration/*} & 支持学生讲解评分、协同摘要和导游修正内容提交。 \\
\bottomrule
\end{longtable}

\section{概要设计}

\subsection{系统总体架构}

系统采用五层架构：客户端层、业务服务层、AI 能力层、知识增强层、数据与运维层。客户端只负责交互和展示，不直接访问模型和数据库；业务服务层负责权限、流程、状态和接口编排；AI 能力层提供模型适配；知识增强层提供可信事实依据；数据与运维层负责持久化、缓存、文件和监控。

\diagramimage[1\textwidth]{design-layer-architecture.png}{系统设计五层架构图}
\newpage
\begin{longtable}{p{2.7cm}p{4.1cm}p{6.2cm}}
\toprule
架构层 & 组成 & 设计说明 \\
\midrule
客户端层 & 游客端、学生端、导游端、管理端 & 负责多模态输入采集、业务页面展示、结果反馈和用户操作提交。游客端优先适配移动场景，管理端和学生端优先适配 PC Web。 \\
业务服务层 & 导览、实训、协同、审核、日志、用户权限 & 负责业务规则、流程编排、权限校验、状态流转和统一 API 输出，是系统核心控制层。 \\
AI 能力层 & ASR、翻译、视觉、LLM、向量化、TTS & 通过 AI 编排服务统一调用模型能力，模型以适配器形式接入，保证业务服务不依赖具体模型实现。 \\
知识增强层 & RAG 知识库、术语表、文化知识图谱 & 为大模型生成提供受控上下文，保证文旅事实、术语翻译和文化关系可追踪。 \\
数据与运维层 & MySQL、PostgreSQL/pgvector 预留、Neo4j、Redis、MinIO、运行日志 & 保存结构化数据、向量索引、图谱关系、缓存、文件对象、审计日志和监控指标。 \\
\bottomrule
\end{longtable}

\subsection{总体技术选型}

\begin{longtable}{p{2.6cm}p{3.8cm}p{6.6cm}}
\toprule
类别 & 技术选型 & 设计理由 \\
\midrule
后端框架 & FastAPI & 支持异步接口、自动生成 OpenAPI 文档，适合快速构建 AI 编排和业务服务接口。 \\
前端框架 & React 18、TypeScript、Vite & 单一 React Web App 覆盖游客、学生、导游、知识库和系统管理入口，游客端优先适配移动浏览器。 \\
业务数据库 & MySQL 8.4 & 保存用户、会话、消息、审核、训练、反馈、日志和管理配置等结构化业务数据。 \\
知识检索 & 应用层哈希向量检索 & 当前使用 256 维哈希向量与余弦相似度完成 RAG 检索；PostgreSQL/pgvector 已启动，作为后续增强预留。 \\
图数据库 & Neo4j & 保存景点、城市、民族、节庆、非遗、饮食、建筑和禁忌等文化实体关系。 \\
缓存 & Redis & 缓存会话上下文、热点问答、临时任务状态、短期检索结果和访问令牌。 \\
文件存储 & MinIO & 保存游客图片、语音文件、学生训练录音、知识文档原文和 TTS 音频。 \\
模型推理 & Ollama、faster-whisper、Windows SAPI & Ollama 提供文本与视觉模型，faster-whisper 提供 ASR，Windows SAPI 提供 TTS；文本模型也支持 OpenAI-compatible Provider。 \\
部署方式 & Docker Compose、本机进程 & Docker Compose 启动数据基础设施；PowerShell 脚本启动 FastAPI、Vite 和本机 Ollama。Kubernetes 属于后续扩展方向。 \\
监控日志 & 健康检查、架构审计、请求追踪 & 当前由 FastAPI 接口提供观测数据，外接监控套件属于后续扩展方向。 \\
\bottomrule
\end{longtable}

\subsection{子系统划分}

\begin{longtable}{p{2.6cm}p{4.2cm}p{6.2cm}}
\toprule
子系统 & 核心模块 & 主要职责 \\
\midrule
用户与权限子系统 & 用户账号、角色权限、登录日志、鉴权中间件 & 管理游客、学生、导游和管理员身份，控制不同角色可访问的页面、接口和数据范围。 \\
智能导览子系统 & 导览会话、消息记录、导览服务、路线方案、反馈 & 处理文本、语音、图片和位置输入，生成可信多语回答、路线建议和文化提示。 \\
知识库管理子系统 & 文档、切片、术语、图谱、审核任务 & 完成文旅资料上传、清洗、切分、审核、向量化、术语维护和文化关系维护。 \\
AI 导游实训子系统 & 训练场景、虚拟游客、训练记录、语音提交、训练报告 & 支持学生模拟带团训练、语音讲解提交、多维评分和报告生成。 \\
真人导游协同子系统 & 会话摘要、接管记录、导游修正、优质案例 & 支持真人导游查看游客上下文、人工接管、修正 AI 回答和沉淀案例。 \\
日志与反馈子系统 & 模型调用日志、检索日志、用户反馈、异常日志 & 支撑问题追踪、内容修订、模型效果分析和运维排障。 \\
\bottomrule
\end{longtable}

\subsection{系统软件结构设计}

系统软件结构按客户端、接口、业务服务、AI 适配、数据访问和基础设施进行分包。客户端包只保存界面、状态管理和 API 调用逻辑；接口包负责请求鉴权、参数校验和路由分发；业务服务包承载导览、实训、协同、审核和日志等核心流程；AI 包通过适配器屏蔽不同模型实现差异；仓储包统一封装数据库、向量库、图数据库、缓存和对象存储访问。

\diagramimage[0.92\textwidth]{design-package-structure.png}{系统软件包结构图}

\begin{longtable}{p{3.2cm}p{4.1cm}p{5.6cm}}
\toprule
软件包 & 主要组成 & 设计说明 \\
\midrule
\code{client} & 游客端、学生端、导游端、管理端 & 负责页面展示、表单校验、多模态输入采集和接口调用，不直接访问数据库和模型服务。 \\
\code{api} & 鉴权中间件、参数校验器、路由控制器 & 负责接收 HTTP 请求，统一处理身份校验、错误响应、限流和接口路由。 \\
\code{services} & 导览、实训、知识、协同、审核、日志服务 & 系统核心业务逻辑所在层，负责流程编排、状态转换和跨模块协作。 \\
\code{ai} & ASR、翻译、视觉、向量化、LLM、TTS 适配器 & 将模型调用封装为稳定接口，便于替换本地模型或云端模型。 \\
\code{repositories} & 用户、会话、知识、训练、审核仓储 & 封装数据读写细节，避免业务服务直接拼接 SQL 或依赖具体存储组件。 \\
\code{infrastructure} & MySQL、PostgreSQL/pgvector、Neo4j、Redis、MinIO & 提供运行时持久化、预留向量能力、缓存、对象存储和图谱镜像。 \\
\bottomrule
\end{longtable}

\subsection{部署概要设计}

当前 MVP 采用“Docker 基础设施 + 本机应用进程”部署方式。Docker Compose 启动 MySQL、PostgreSQL/pgvector、Redis、MinIO 和 Neo4j；启动脚本校验 Ollama 文本与视觉模型后，以本机进程启动 FastAPI 和 Vite。导览、实训、协同、知识和日志能力在 FastAPI 应用内按模块组织；独立服务拆分与 Kubernetes 编排属于后续扩展方向。

\begin{longtable}{p{3cm}p{4cm}p{6cm}}
\toprule
节点/容器 & 部署内容 & 说明 \\
\midrule
Web 进程 & React Vite Web App & 提供游客、学生、导游、知识库和系统管理页面，开发环境由 Vite 提供访问入口。 \\
API 容器 & FastAPI 应用、鉴权中间件、业务服务 & 对客户端提供 HTTP API，对内部服务提供模块接口。 \\
本机模型服务 & Ollama 文本模型、Ollama 视觉模型、faster-whisper、Windows SAPI & 由 Provider 和适配器调用，可通过环境变量替换文本模型 Provider。 \\
数据容器 & MySQL、PostgreSQL/pgvector、Neo4j、Redis、MinIO & MySQL 保存运行时数据；Neo4j、Redis、MinIO 已接入；PostgreSQL/pgvector 为向量增强预留。 \\
运行观测 & 健康检查、系统指标、模型日志、请求追踪 & 当前由 FastAPI 接口提供观测数据，外接监控套件作为后续增强。 \\
\bottomrule
\end{longtable}

\diagramimage[0.7\textwidth]{design-deployment-topology.png}{部署拓扑设计图}

\subsection{关键业务流程概要}

\subsubsection{智能导览流程}

智能导览流程统一处理文本、语音和图片输入。语音先经过 ASR 转写，图片先经过视觉模型和 OCR 识别，随后进入标准文本问题处理流程。系统将问题翻译或标准化为统一语言，检索审核通过的知识切片，补充图谱关系，再由大语言模型基于受控上下文生成回答。

\flowdiagramimage[0.5\textwidth]{design-guide-orchestration-activity.png}{智能导览编排活动图}
\newpage
\subsubsection{知识库管理流程}

知识库管理流程从资料上传开始，经过格式解析、文本清洗、知识切片、元数据标注和人工审核。审核通过的切片才会发布到正式知识库；当前应用层检索直接读取已发布知识，向量状态字段为后续增强保留；退回或下线内容不能进入游客端回答。

\flowdiagramimage[0.62\textwidth]{design-knowledge-review-activity.png}{知识入库审核活动图}
\newpage
\subsubsection{AI 导游实训流程}

实训流程以训练场景为入口，系统生成虚拟游客问题，学生提交语音讲解后由 ASR 转写，并结合知识库标准内容和评分规则生成报告。评分结果作为教学辅助，不替代教师最终评价。

\flowdiagramimage[0.35\textwidth]{design-training-evaluation-activity.png}{AI 导游实训评分活动图}
\newpage
\subsubsection{真人导游协同流程}

导游协同流程连接 AI 服务和人工服务。系统在游客会话中持续记录问题、语言、地点、兴趣主题和 AI 回答摘要，导游可查看摘要并接管复杂会话。导游修正和优质案例进入审核流程，通过后回流知识库或案例库。

\flowdiagramimage[0.7\textwidth]{design-human-guide-collaboration-activity.png}{真人导游协同活动图}
\newpage
\section{详细设计}

\subsection{总体类结构设计}

系统总体类结构由用户、会话、消息、知识文档、知识切片、训练记录、导游修正和模型调用日志等核心实体组成。用户与导览会话、训练记录、导游修正存在一对多关系；会话包含多条消息；知识文档切分为多个知识切片；消息产生模型调用日志和导游修正记录。

\diagramimage[0.84\textwidth]{design-core-domain-class.png}{系统核心领域类图}

\subsection{用户与权限子系统详细设计}

\subsubsection{类图}

\diagramimage[0.82\textwidth]{design-user-auth-class.png}{用户与权限设计类图}

\subsubsection{类职责}

\begin{longtable}{p{3cm}p{4.2cm}p{6cm}}
\toprule
类名 & 主要属性/方法 & 设计说明 \\
\midrule
User & userId、username、roleId、status、login()、disable() & 保存账号信息和账号状态，是权限判断的主体对象。 \\
Role & roleId、roleName、addPermission()、removePermission() & 描述游客、学生、导游、管理员等角色，并维护角色权限集合。 \\
Permission & permissionCode、resource、action、match() & 描述可访问资源和操作动作，用于接口级和页面级权限控制。 \\
LoginLog & loginTime、ipAddress、result、recordLogin() & 记录登录行为，支持安全审计和异常访问追踪。 \\
AuthService & authenticate()、checkPermission()、assignRole() & 完成认证、授权、角色分配和权限校验。 \\
\bottomrule
\end{longtable}

\subsubsection{权限矩阵}

\begin{longtable}{p{2.5cm}p{2.4cm}p{2.4cm}p{2.4cm}p{2.4cm}}
\toprule
功能资源 & 游客 & 学生 & 导游 & 管理员 \\
\midrule
智能导览问答 & 可用 & 可用 & 可查看 & 可查看日志 \\
图片/语音问答 & 可用 & 可用 & 可查看摘要 & 可查看日志 \\
AI 实训 & 不可用 & 可用 & 可查看案例 & 可配置场景 \\
会话接管 & 不可用 & 不可用 & 可用 & 可查看记录 \\
知识文档上传 & 不可用 & 不可用 & 可提交案例 & 可用 \\
内容审核 & 不可用 & 不可用 & 不可用 & 可用 \\
模型调用日志 & 不可用 & 不可用 & 仅关联会话 & 可用 \\
\bottomrule
\end{longtable}

\subsection{智能导览子系统详细设计}

\subsubsection{类图}

\diagramimage[1\textwidth]{design-guide-service-class.png}{智能导览服务类图}
\newpage
\subsubsection{模块职责}

\begin{longtable}{p{3.2cm}p{4.2cm}p{5.8cm}}
\toprule
模块 & 输入/输出 & 设计说明 \\
\midrule
导览会话管理 & 输入用户、景区、语种；输出 sessionId & 保存连续对话上下文，支持游客和导游围绕同一会话协作。 \\
多模态输入处理 & 输入文本、音频、图片、位置；输出标准 query & 将语音和图片转换为可检索文本，统一进入 RAG 问答流程。 \\
RAG 检索调用 & 输入 query、spotId、language；输出知识片段列表 & 仅召回审核通过、适用景区和语言范围匹配的知识切片。 \\
回答生成 & 输入 query、知识片段、图谱关系；输出回答文本 & 调用 LLM 生成导览回答，必须约束模型只基于上下文回答事实。 \\
结果输出 & 输入回答、来源、目标语言；输出前端响应 & 进行术语校验、翻译和可选 TTS 合成，并保存消息记录。 \\
\bottomrule
\end{longtable}

\subsubsection{核心处理规则}

\begin{enumerate}
  \item 文本问题直接进入语种识别和意图识别流程。
  \item 语音问题先保存原始音频，再调用 ASR 转写，并将转写文本送入文本问答流程。
  \item 图片问题先调用视觉模型识别对象、OCR 文本和地点线索，再生成检索关键词。
  \item 涉及文化事实的问题必须检索知识库；缺少可信知识时返回“暂无可靠资料”。
  \item 所有回答必须保存 question、answer、sourceChunks、modelVersion、latency 和 status。
\end{enumerate}
\newpage
\subsection{知识库管理子系统详细设计}

\subsubsection{类图}

\diagramimage[1\textwidth]{design-knowledge-service-class.png}{知识库管理服务类图}

\subsubsection{知识入库状态机}

\begin{longtable}{p{2.8cm}p{4cm}p{6.4cm}}
\toprule
状态 & 触发动作 & 说明 \\
\midrule
\code{UPLOADED} & 管理员上传文档 & 原始文件保存到 MinIO，元数据写入 MySQL。 \\
\code{PARSED} & 文档解析完成 & 文档完成格式解析、文本抽取和基础清洗。 \\
\code{CHUNKED} & 知识切片完成 & 文本按主题、段落、问答对或景区维度切分。 \\
\code{PENDING_REVIEW} & 创建审核任务 & 切片等待管理员审核，暂不可进入游客端回答。 \\
\code{APPROVED} & 审核通过 & 切片可向量化并发布到正式知识库。 \\
\code{REJECTED} & 审核退回 & 内容需修改或补充来源，不进入检索。 \\
\code{PUBLISHED} & 向量化入库完成 & MySQL 保存正文和向量状态，正式知识通过应用层检索；pgvector 接入属于后续增强。 \\
\code{OFFLINE} & 管理员下线 & 内容失效或存在风险，检索时必须过滤。 \\
\bottomrule
\end{longtable}
\newpage
\subsubsection{知识切片元数据}

\begin{longtable}{p{3cm}p{3.4cm}p{6.6cm}}
\toprule
字段 & 类型 & 说明 \\
\midrule
chunkId & String & 知识切片唯一编号。 \\
documentId & String & 来源文档编号。 \\
spotId & String & 适用景区或地点编号。 \\
language & String & 原文语种。 \\
topicTags & String/List & 景点、民族、节庆、非遗、饮食、路线等主题标签。 \\
content & Text & 切片正文。 \\
sourceName & String & 来源名称，如文旅局资料、景区资料、课程资料等。 \\
reviewStatus & Enum & 审核状态，正式检索只允许 APPROVED/PUBLISHED。 \\
sensitiveLevel & Integer & 敏感级别，用于民族、宗教、政策等内容过滤。 \\
updatedAt & DateTime & 内容更新时间，用于动态信息时效判断。 \\
\bottomrule
\end{longtable}

\subsection{AI 导游实训子系统详细设计}

\subsubsection{类图}

\diagramimage[1\textwidth]{design-training-service-class.png}{AI 导游实训服务类图}
\newpage
\subsubsection{训练场景设计}

\begin{longtable}{p{2.8cm}p{4cm}p{6.3cm}}
\toprule
场景类型 & 示例任务 & 评分重点 \\
\midrule
景点讲解 & 介绍大理古城、丽江古城、石林等景点 & 内容准确性、结构完整度、讲解逻辑。 \\
民族文化讲解 & 解释白族、傣族、纳西族、彝族文化 & 文化准确性、礼仪敏感性、术语统一。 \\
接团服务 & 到达集合、行程说明、注意事项提醒 & 服务礼仪、表达清晰度、流程完整度。 \\
投诉处理 & 路线变更、等待时间、服务争议处理 & 应变能力、情绪安抚、解决方案合理性。 \\
跨文化交流 & 面向不同国家游客解释文化差异 & 多语表达、跨文化敏感性、禁忌规避。 \\
应急处理 & 迷路、天气变化、游客身体不适 & 风险识别、响应及时性、安全提示。 \\
\bottomrule
\end{longtable}

\subsubsection{评分维度}

\begin{longtable}{p{2.8cm}p{2cm}p{7.8cm}}
\toprule
维度 & 建议权重 & 评价说明 \\
\midrule
内容准确性 & 30\% & 讲解内容是否符合审核通过的知识库材料，是否存在事实错误。 \\
讲解完整度 & 20\% & 是否覆盖景点背景、文化关系、注意事项和游客可能关心的问题。 \\
表达清晰度 & 15\% & 语言是否流畅，逻辑是否清楚，是否有明显停顿或歧义。 \\
服务礼仪 & 10\% & 是否符合导游服务规范，是否包含必要提醒。 \\
应变能力 & 10\% & 面对追问、投诉和突发情况时是否给出合理处理。 \\
跨文化敏感性 & 10\% & 是否避免民族、宗教、礼仪禁忌方面的不当表达。 \\
多语能力 & 5\% & 外语或多语表达是否准确、自然，术语是否一致。 \\
\bottomrule
\end{longtable}
\newpage
\subsubsection{实训评分增强设计}

Topo-SCTB-Net 可作为 AI 导游实训评分模块的后续研究型增强方案，当前仓库未接入该模型。该模块面向中文导游实训评价，输入学生讲解文本、题目、场景、游客背景、SNEU 场景化规范证据单元、知识库、服务规范和评分 rubric，输出事实/关系状态、规范状态、受众适配状态、证据 span、风险 span、知识分、表达分、服务分和最终分。

\begin{longtable}{p{3.2cm}p{4.2cm}p{5.8cm}}
\toprule
组件 & 输入/输出 & 设计说明 \\
\midrule
SNEU 证据单元 & 输入场景规范、知识来源、评分规则；输出结构化证据单元 & 将导游讲解需要满足的事实、关系、礼仪和受众条件结构化。 \\
Soft Evidence-State Graph & 输入候选 span、SNEU slot、游客条件；输出软边图 & 表达学生回答、知识证据和状态判断之间的关系。 \\
Topology-guided CTB Loss & 输入反事实干预和状态变化；输出训练约束 & 约束目标状态翻转，同时保持非目标状态稳定，提高评分鲁棒性。 \\
Monotonic Scorer & 输入 fact、relation、norm、audience 状态；输出分项和总分 & 保证状态效用单调，使 supported 优于 partial，violated 低于 risky。 \\
Evidence-grounded Feedback & 输入证据 span 和风险 span；输出反馈建议 & 给出基于证据的模板化反馈，减少纯 LLM 评价的不稳定性。 \\
\bottomrule
\end{longtable}

实训评分由 \code{TrainingEvaluator} 服务统一完成。该服务综合知识库召回、事实覆盖度、礼仪关键词、讲解结构、多语表达质量，并可在后续接入 Topo-SCTB-Net 证据评分结果。当前评分请求由 FastAPI 同步处理并直接返回训练报告。
\newpage
\subsection{真人导游协同子系统详细设计}

\subsubsection{类图}

\diagramimage[0.82\textwidth]{design-human-guide-service-class.png}{真人导游协同服务类图}

\subsubsection{协同规则}

\begin{enumerate}
  \item 默认情况下游客与 AI 导览系统交互，系统持续生成会话摘要。
  \item 当问题涉及投诉、复杂历史文化、宗教民族敏感内容或 AI 多次无法回答时，可提示导游接管。
  \item 导游接管后，游客端应显示人工服务状态，避免游客误以为仍由 AI 自动回答。
  \item 导游修正内容、优质讲解案例和人工补充资料必须进入审核流程。
  \item 审核通过后内容才能进入知识库、案例库或实训训练材料。
\end{enumerate}

\subsection{接口详细设计}

\subsubsection{客户端 API 接口}

当前前端直接调用 FastAPI 路由，不设置额外的 \code{/api/guide/*} 或 \code{/api/admin/*} 别名前缀。主要接口如下。

\begin{longtable}{p{5.0cm}p{1.35cm}p{3.15cm}p{3.0cm}}
\toprule
接口路径 & 方法 & 请求参数 & 响应内容 \\
\midrule
\code{/api/auth/login} & POST & username、password & token、user \\
\code{/api/chat} & POST & question、language、location、session\_id & answer、sources、provider、model、reliable \\
\code{/api/audio/ask} & POST & file、language、location、session\_id & transcript、answer、sources \\
\code{/api/image/ask} & POST & file、question、language、location、session\_id & recognition、answer、sources \\
\code{/api/route/recommend} & POST & visitor\_type、interest、language & route、relatedItems \\
\code{/api/training/scenarios} & GET & 无 & items \\
\code{/api/training/score} & POST & scenario、question、answer & total、metrics、feedback、judgeMode \\
\code{/api/training/records} & GET & 无 & items \\
\code{/api/collaboration/sessions} & GET & 无 & items \\
\code{/api/sessions/{session\_id}/takeover} & POST & note & sessionStatus \\
\code{/api/sessions/{session\_id}/guide-reply} & POST & answer & message \\
\code{/api/sessions/{session\_id}/release} & POST & 无 & sessionStatus \\
\code{/api/collaboration/correction} & POST & question、ai\_answer、guide\_note & correction、reviewTask \\
\code{/api/knowledge/documents} & POST & file、source、tags & documentId、chunkCount、tasks \\
\code{/api/review/tasks} & GET & 无 & items \\
\code{/api/review/tasks/{task\_id}/decision} & POST & status、reviewer、comment & reviewTask \\
\code{/api/logs/model-calls} & GET & 无 & items \\
\code{/api/logs/request-traces} & GET & 无 & items \\
\bottomrule
\end{longtable}
\newpage
\subsubsection{接口扩展说明}

知识切片、术语、图谱、用户、角色权限、系统设置、系统指标、游客偏好、收藏、反馈和导游案例均已提供资源化 CRUD 或查询接口。接口的权威清单以 FastAPI 自动生成的 \code{/docs} 页面和 \code{backend/app/main.py} 为准；前端调用路径与后端路由保持一致。

\subsubsection{界面设计}

系统界面按角色划分为游客导览端、学生实训端、导游协同端、管理端和日志监控端。界面设计遵循“角色入口清晰、主任务优先、结果可追溯、异常可解释”的原则：移动端优先突出快速提问、语音输入、图片识别和路线推荐；学生端突出训练任务、讲解提交和评分反馈；导游端突出待接管会话、风险摘要和修正提交；管理端突出知识入库、审核发布、术语维护、图谱维护和日志追踪。

\begin{longtable}{p{2.4cm}p{3.1cm}p{3.6cm}p{3.9cm}}
\toprule
角色端 & 界面名称 & 核心区域 & 设计说明 \\
\midrule
游客端 & 导览首页 & 语言选择、文本提问、语音按钮、拍照入口、路线推荐入口 & 面向景区移动场景，优先突出快速提问、多语回答和来源提示。 \\
游客端 & 多语问答结果页 & 回答卡片、来源列表、可靠性提示、语音播报、继续追问 & 展示答案生成依据，避免无来源文化事实直接呈现。 \\
游客端 & 图片识别问答页 & 图片预览、识别摘要、补充问题、知识检索来源、回答结果 & 图片信息只作为输入线索，最终回答仍需进入知识检索链路。 \\
游客端 & 路线推荐页 & 当前位置、兴趣标签、游览时长、路线节点、注意事项 & 支持按时间、兴趣、语言和人群类型生成路线。 \\
学生端 & 实训场景页 & 场景列表、难度标签、虚拟游客设定、任务目标 & 使学生在提交前明确讲解对象、语言要求和评分重点。 \\
学生端 & 讲解提交页 & 题目内容、录音控件、文本稿、提交按钮、重录入口 & 支持语音和文本两类提交方式，便于不同训练环境使用。 \\
学生端 & 评分报告页 & 总分、分项得分、事实覆盖、语言表达、改进建议 & 将评分结果拆解到可训练的能力维度。 \\
学生端 & 训练历史页 & 历史记录、场景筛选、得分趋势、报告入口 & 支持学生回看进步趋势和薄弱维度。 \\
导游端 & 协同会话列表页 & 风险等级、游客语言、最新问题、等待时长、接管入口 & 帮助真人导游优先处理高风险或高等待时长会话。 \\
导游端 & 会话接管页 & 会话摘要、原始问答、建议回复、人工输入、发送按钮 & 接管后保留 AI 上下文，减少导游重复询问成本。 \\
导游端 & 修正提交页 & 原回答、修正内容、依据说明、关联景点、提交审核 & 将人工修正转为可审核知识资产。 \\
管理端 & 知识文档管理页 & 上传区域、文档列表、版本号、来源、状态筛选 & 支撑知识入库、版本追踪和来源管理。 \\
管理端 & 知识审核页 & 切片详情、相似内容、敏感等级、审核意见、通过/退回/下线 & 保证正式检索只使用审核通过内容。 \\
管理端 & 术语维护页 & 术语列表、语种字段、标准译名、适用范围、冲突提示 & 统一重点文旅名词、多语译名和敏感表达。 \\
管理端 & 文化图谱维护页 & 实体列表、关系编辑、图谱预览、影响范围 & 维护景点、民族、节庆、饮食、禁忌等文化关系。 \\
管理端 & 用户权限页 & 用户列表、角色分配、权限矩阵、登录日志 & 支持按游客、学生、导游、管理员进行权限隔离。 \\
监控端 & 模型调用日志页 & 请求编号、模型名称、耗时、输入摘要、错误信息 & 用于排查模型异常、响应超时和质量问题。 \\
监控端 & 运行状态页 & 服务健康、数据库状态、队列长度、告警列表 & 展示系统运行状态和架构审计结果。 \\
\bottomrule
\end{longtable}

\paragraph{游客端界面截图}

\systemscreenimage[0.86\textwidth]{tourist-home.png}{游客导览首页}

\systemscreenimage[0.86\textwidth]{tourist-chat.png}{游客多语问答页}

\systemscreenimage[0.86\textwidth]{tourist-image.png}{游客图片识别问答页}

\systemscreenimage[0.86\textwidth]{tourist-routes.png}{游客路线推荐页}

\paragraph{学生端界面截图}

\systemscreenimage[0.86\textwidth]{student-scenarios.png}{学生实训场景页}

\systemscreenimage[0.86\textwidth]{student-training.png}{学生讲解提交页}

\systemscreenimage[0.86\textwidth]{student-reports.png}{学生评分报告页}

\systemscreenimage[0.86\textwidth]{student-reports.png}{学生训练历史页}

\paragraph{导游端界面截图}

\systemscreenimage[0.86\textwidth]{guide-sessions.png}{导游协同会话列表页}

\systemscreenimage[0.86\textwidth]{guide-takeover.png}{导游会话接管页}

\systemscreenimage[0.86\textwidth]{guide-corrections.png}{导游修正提交页}

\paragraph{管理端界面截图}

\systemscreenimage[0.9\textwidth]{knowledge-documents.png}{管理端知识文档管理页}

\systemscreenimage[0.9\textwidth]{knowledge-review.png}{管理端知识审核页}

\systemscreenimage[0.9\textwidth]{knowledge-terms.png}{管理端术语维护页}

\systemscreenimage[0.9\textwidth]{knowledge-graph.png}{管理端文化图谱维护页}

\systemscreenimage[0.9\textwidth]{system-users.png}{管理端用户权限页}

\paragraph{监控端界面截图}

\systemscreenimage[0.9\textwidth]{system-logs.png}{模型调用日志页}

\systemscreenimage[0.9\textwidth]{system-health.png}{系统运行状态页}

\subsubsection{内部服务接口}

\begin{longtable}{p{5.4cm}p{3.7cm}p{3.5cm}}
\toprule
接口 & 输入 & 输出 \\
\midrule
\code{AuthService.checkPermission} & userId、resource、action & allowed、denyReason \\
\code{GuideService.submitQuestion} & sessionId、query、inputType & messageId、answer、sources \\
\code{AIOrchestrator.process} & query、inputType、targetLanguage & normalizedQuery、answer、trace \\
\code{KnowledgeService.search} & query、spotId、language、topK & reviewedChunks、scores \\
\code{GraphService.queryRelations} & entityNames、relationTypes & relationResult \\
\code{TermService.applyTerms} & text、targetLanguage、topicTags & normalizedText、termHits \\
\code{ReviewService.createTask} & objectType、objectId、riskLevel & taskId、status \\
\code{LogService.recordModelCall} & requestId、modelName、input、output、latency & logId \\
\code{TrainingService.scoreSubmission} & recordId、transcript、sceneId & scoreDetail、reportId \\
\bottomrule
\end{longtable}

\subsubsection{模型适配接口}

\begin{longtable}{p{3.7cm}p{3.2cm}p{6.0cm}}
\toprule
适配器 & 推荐模型 & 接口职责 \\
\midrule
ASRAdapter & Whisper / faster-whisper & 输入音频文件，输出转写文本、语种和置信度。 \\
TranslationAdapter & 术语表替换适配器 & 当前优先应用术语表标准译名；可后续接入外部翻译模型。 \\
VisionAdapter & Ollama qwen3-vl:4b & 识别图片摘要，并将摘要送入可信知识检索链路。 \\
EmbeddingAdapter & 哈希向量 & 当前生成 256 维哈希向量并计算余弦相似度；pgvector 作为后续增强预留。 \\
LLMAdapter & Ollama qwen3.5:9b / OpenAI-compatible & 基于检索上下文生成问答；Provider 可通过环境变量切换。 \\
TTSAdapter & Windows SAPI & 将回答文本转换为可播放 WAV 音频。 \\
\bottomrule
\end{longtable}

\subsection{数据存储详细设计}

\subsubsection{核心 ER 关系}

\diagramimage[0.68\textwidth]{design-data-storage-er.png}{数据存储 ER 设计图}

\subsubsection{数据库表设计}

当前物理数据库为 MySQL 8.4。后端在启动时由 \code{runtime\_store.py} 创建运行时表，CSV 数据由 \code{bootstrap\_mysql.py} 导入种子内容表。主键使用字符串形式的 UUID，JSON 列表与链路详情以 JSON 文本保存，时间字段使用 ISO 8601 字符串。

\begin{longtable}{p{3.2cm}p{4.1cm}p{5.7cm}}
\toprule
表名 & 核心字段 & 当前用途 \\
\midrule
\code{users} & id、username、password、role、language、status & 用户认证与角色入口。 \\
\code{guide\_sessions} & id、visitor\_name、language、location、status、risk\_level & 游客会话与人工接管状态。 \\
\code{messages} & session\_id、role、input\_type、question、answer、sources\_json、reliable & 保存问答、图片、语音和导游回复。 \\
\code{visitor\_profiles} & session\_id、interests\_json、intent\_types\_json、visit\_status、risk\_level & 保存会话画像。 \\
\code{request\_traces} & session\_id、question、language、pipeline\_json & 保存 RAG 检索和问答链路。 \\
\code{model\_call\_logs} & request\_id、capability、provider、model、latency\_ms、status & 保存模型调用日志。 \\
\code{knowledge\_documents} & title、source、tags\_json、content、status、vector\_status & 保存上传文档。 \\
\code{knowledge\_chunks} & document\_id、title、content、tags\_json、status、vector\_status & 保存切片和向量状态。 \\
\code{review\_tasks} & object\_type、object\_id、title、content、status、reviewer、comment & 保存知识与导游修正审核任务。 \\
\code{training\_records} & scenario、question、answer、score、metrics\_json、feedback\_json、judge\_mode & 保存实训评分报告。 \\
\code{guide\_corrections} & record\_id、mode、guide\_note、optimized\_answer、status & 保存导游修正。 \\
\code{feedback} & message\_id、rating、content、status & 保存用户反馈。 \\
\code{favorites}、\code{takeover\_logs} & session\_id、业务对象、动作、备注 & 保存游客收藏与人工接管历史。 \\
\code{roles}、\code{permissions}、\code{role\_permissions} & 角色、权限和关联编号 & 保存后台权限矩阵。 \\
\code{system\_settings} & values\_json、updated\_at & 保存系统设置。 \\
\bottomrule
\end{longtable}

\paragraph{种子内容表}
\code{knowledge\_items}、\code{graph\_relations}、\code{routes}、\code{route\_filters}、\code{training\_scenarios}、\code{guide\_places}、\code{guide\_questions} 和 \code{collaboration\_cases} 由 CSV 初始化并通过 \code{store.py} 读写。

\paragraph{扩展说明}
PostgreSQL/pgvector 已随 Docker Compose 启动，但当前正式 RAG 检索仍由应用层哈希向量完成。后续接入 pgvector 时，应保持审核状态过滤规则不变。

\subsubsection{非关系数据设计}

\begin{longtable}{p{2.8cm}p{3.7cm}p{6.7cm}}
\toprule
存储组件 & 数据内容 & 设计说明 \\
\midrule
应用层哈希向量检索 & title、content、tags、score & 当前使用哈希向量和余弦相似度检索 MySQL 中已发布知识；PostgreSQL/pgvector 作为后续扩展。 \\
Neo4j & 景点、城市、民族、节庆、非遗、建筑、饮食、禁忌、路线关系 & 支持文化关联解释和路线推荐补充。 \\
Redis & sessionContext、hotQuestions、token、temporaryTask & 缓存短期会话和热点结果，降低数据库压力。 \\
MinIO & audio、image、document、ttsAudio & 保存二进制文件，数据库仅保存对象地址和元数据。 \\
MySQL 日志表 & requestTraces、modelCallLogs、feedback & 支撑模型调用追踪、异常排查和反馈闭环。 \\
\bottomrule
\end{longtable}

\subsection{异常与安全设计}

\subsubsection{异常处理表}

\begin{longtable}{p{3.3cm}p{4.2cm}p{5.4cm}}
\toprule
异常场景 & 前端提示 & 后端处理 \\
\midrule
未登录访问受限功能 & 请先登录后继续操作 & 返回 401，记录访问资源和来源 IP。 \\
角色权限不足 & 当前账号无权访问该功能 & 返回 403，记录 userId、resource、action。 \\
知识库无可信资料 & 暂无可靠资料 & 禁止 LLM 自由编造，记录 missedQuery 用于知识补充。 \\
ASR 失败 & 语音识别失败，请重新录制 & 保留音频对象和错误信息，允许重试。 \\
图片识别失败 & 图片内容暂无法识别 & 记录 VisionAdapter 错误，允许改用文字提问。 \\
模型服务超时 & AI 服务繁忙，请稍后再试 & 触发重试或降级，记录模型、耗时和 traceId。 \\
审核不通过 & 内容未通过审核 & 保存审核意见，禁止发布到正式知识库。 \\
文件上传失败 & 文件格式或大小不符合要求 & 校验扩展名、MIME、大小和病毒扫描结果。 \\
\bottomrule
\end{longtable}

\subsubsection{安全控制}

\begin{enumerate}
  \item 接口认证使用访问令牌或服务端会话，管理端、导游端和学生端必须登录。
  \item 角色权限按资源和动作控制，避免学生访问管理端、导游访问全量模型日志。
  \item 用户输入、上传文件和知识文档均需格式校验，图片和音频文件不直接执行。
  \item 敏感内容、民族宗教内容和政策类内容必须经过审核状态过滤。
  \item 日志中应对手机号、证件号、定位细节等敏感字段脱敏。
  \item 管理员审核、导游修正、知识发布和权限变更必须记录审计日志。
\end{enumerate}

\subsection{可观测性与维护设计}

\begin{longtable}{p{2.8cm}p{4.2cm}p{6.2cm}}
\toprule
对象 & 指标/日志 & 维护用途 \\
\midrule
HTTP API & 请求量、状态码、耗时、错误率 & 发现接口异常、性能瓶颈和高频访问路径。 \\
模型调用 & 模型名称、版本、输入长度、耗时、失败原因 & 判断模型服务是否稳定，支持模型版本回溯。 \\
RAG 检索 & query、召回 chunk、相似度、过滤原因 & 分析知识库覆盖率和检索质量。 \\
审核流程 & taskId、对象类型、处理人、处理意见、状态变化 & 支撑内容治理和责任追踪。 \\
用户反馈 & messageId、评分、反馈内容、处理状态 & 形成知识修订和模型优化闭环。 \\
基础设施 & CPU、内存、GPU、磁盘、网络、容器状态 & 支撑运维告警、扩容和故障定位。 \\
\bottomrule
\end{longtable}

\subsection{设计约束与运行保障}

\begin{enumerate}
  \item 新增语种时，应优先配置术语表、翻译适配器和 TTS 参数，不修改导览主流程。
  \item 新增景区时，应通过知识文档、知识切片、图谱实体和景点配置入库完成。
  \item 新增模型时，只实现对应 Adapter，并在 AI 编排配置中切换模型名称和版本。
  \item 新增训练场景时，只新增场景模板、评分规则和虚拟游客配置。
  \item 当并发量提升时，优先拆分模型推理服务、导览服务和知识检索服务，并对热点查询启用 Redis 缓存。
\end{enumerate}
